\item Una persona que busca en un recipiente vacío es capaz de ver el extremo más alejado de la parte inferior del contenedor (figura P34.20a). La altura del contenedor es $h$ y su anchura es $d$. Cuando el recipiente está completamente lleno de un fluido con índice de refracción $n$ y se ve bajo el mismo ángulo, la persona puede ver el centro de una moneda en el centro de la parte inferior del recipiente.
\begin{enumerate}
    \item Demostrar que la razón $h/d$ está dada por:
    $$ \frac{h}{d} = \sqrt{\frac{n^2 - 1}{4 - n^2}} $$
    \item Suponga que el contenedor tiene una anchura de $8.00\text{ cm}$ y está lleno de agua. Use la expresión anterior para encontrar la altura del recipiente.
    \item ¿Para qué rango de valores de $n$ la moneda no será visible para cualquier valor de $h$ y $d$?
\end{enumerate}